\documentclass[12pt,a4paper]{amsart}
\usepackage{amsmath,amssymb,amsthm}
\usepackage{mathtools}
\usepackage[utf8]{inputenc}
\usepackage[T1]{fontenc}
\usepackage{hyperref}
\usepackage{enumitem,booktabs,array}
\usepackage{geometry}
\geometry{margin=2.5cm}

\newtheorem{theorem}{Theorem}[section]
\newtheorem{lemma}[theorem]{Lemma}
\newtheorem{corollary}[theorem]{Corollary}
\newtheorem{proposition}[theorem]{Proposition}
\newtheorem{conjecture}[theorem]{Conjecture}
\theoremstyle{definition}
\newtheorem{definition}[theorem]{Definition}
\newtheorem{example}[theorem]{Example}
\theoremstyle{remark}
\newtheorem{remark}[theorem]{Remark}

\author{Michael Fuhrmann}
\address{Specialist Mathematics Project, Build 141}

\title{Zeta Values at Odd Integers:\\
Irrationality Results, Apéry's Theorem, and Open Problems}

\begin{abstract}
The Riemann zeta function $\zeta(s) = \sum_{n=1}^{\infty} n^{-s}$ evaluated at
positive integers has fascinated mathematicians since Euler's time. While the
values $\zeta(2n)$ at positive even integers are completely understood---they
equal rational multiples of powers of $\pi$---the arithmetic nature of the odd
values $\zeta(2n+1)$ remains largely mysterious. The sole exception is $\zeta(3)$,
proved irrational by Roger Apéry in 1979 in a result that stunned the mathematical
community. We present a comprehensive treatment of zeta values at odd integers,
including Apéry's proof of the irrationality of $\zeta(3)$, the Ball--Rivoal theorem
(2001) guaranteeing infinitely many irrational values among $\{\zeta(2k+1)\}$,
and Zudilin's refinement (2004) that at least one of $\zeta(5), \zeta(7), \zeta(9),
\zeta(11)$ is irrational. We discuss Bernoulli numbers, multiple zeta values,
$p$-adic zeta functions, and connections to modular forms. The irrationality of
$\zeta(5)$ and the transcendence of $\zeta(3)$ remain open conjectures.
\end{abstract}

\maketitle

\tableofcontents

% ============================================================
\section{Introduction}
% ============================================================

The Riemann zeta function is arguably the most important special function in
analytic number theory. Defined for $\operatorname{Re}(s) > 1$ by the absolutely
convergent Dirichlet series
\[
  \zeta(s) = \sum_{n=1}^{\infty} \frac{1}{n^s},
\]
it extends meromorphically to all of $\mathbb{C}$ with a single simple pole at
$s = 1$ (residue~$1$). Euler's product formula
\[
  \zeta(s) = \prod_{p \text{ prime}} \frac{1}{1 - p^{-s}}, \quad
  \operatorname{Re}(s) > 1,
\]
links the zeta function to the distribution of prime numbers, culminating in the
Riemann Hypothesis, which asserts that all non-trivial zeros lie on the critical
line $\operatorname{Re}(s) = \tfrac{1}{2}$.

Beyond its role in prime number theory, the values $\zeta(n)$ at positive integers
have been studied intensively. Euler showed in~1740 that
\[
  \zeta(2n) = (-1)^{n+1} \frac{(2\pi)^{2n} B_{2n}}{2\,(2n)!}, \quad n \geq 1,
\]
where $B_{2n}$ are the Bernoulli numbers. In particular, $\zeta(2) = \pi^2/6$,
$\zeta(4) = \pi^4/90$, and so on---all transcendental numbers, explicitly expressed
in terms of $\pi$.

The situation for \emph{odd} positive integers is strikingly different. No
analogous closed form is known. The question of whether $\zeta(3), \zeta(5),
\zeta(7), \ldots$ are irrational, transcendental, or even algebraically independent
over $\mathbb{Q}(\pi)$ is one of the deepest open problems in number theory.

\subsection*{Historical overview}

\begin{itemize}[leftmargin=2em]
  \item \textbf{Euler (1740)}: Proved $\zeta(2n) \in \mathbb{Q} \cdot \pi^{2n}$.
        The odd case remained open.
  \item \textbf{Apéry (1978/79)}: Proved $\zeta(3) \notin \mathbb{Q}$---a
        \emph{mathematical sensation}, presented at the Journées Arithmétiques in
        Luminy (1978).
  \item \textbf{Beukers (1979)}: Gave an elegant reformulation of Apéry's proof
        using Legendre polynomials and double integrals.
  \item \textbf{Rivoal (2000) and Ball--Rivoal (2001)}: Proved that infinitely many
        odd zeta values are irrational.
  \item \textbf{Zudilin (2004)}: Proved that at least one of $\zeta(5), \zeta(7),
        \zeta(9), \zeta(11)$ is irrational.
\end{itemize}

This paper is organised as follows. Section~2 reviews the fully understood even
case. Section~3 presents Apéry's theorem on $\zeta(3)$ in detail.
Section~4 discusses the constant $\zeta(3)$ and its numerical and historical
significance. Section~5 covers the Ball--Rivoal and Zudilin results.
Section~6 introduces multiple zeta values. Section~7 reviews Bernoulli numbers.
Section~8 covers $p$-adic zeta functions. Section~9 connects the theory to
modular forms. Section~10 summarises open problems. Section~11 gives references.

% ============================================================
\section{Even Values: The Euler Formula}
% ============================================================

\subsection{Statement and proof sketch}

\begin{theorem}[Euler, 1740]
  For every positive integer $n \geq 1$,
  \[
    \zeta(2n) = (-1)^{n+1} \frac{(2\pi)^{2n} B_{2n}}{2\,(2n)!},
  \]
  where $B_{2n}$ denotes the $2n$-th Bernoulli number.
\end{theorem}

\begin{proof}[Proof sketch]
  The classical proof uses the product expansion of $\sin(\pi z)$:
  \[
    \frac{\sin(\pi z)}{\pi z} = \prod_{n=1}^{\infty} \left(1 - \frac{z^2}{n^2}\right).
  \]
  Taking logarithms and differentiating, or expanding both sides as power series
  in $z^2$ and comparing coefficients, yields
  \[
    \sum_{n=1}^{\infty} \frac{1}{n^{2k}} = (-1)^{k+1} \frac{(2\pi)^{2k}
    B_{2k}}{2\,(2k)!}.
  \]
  An alternative route uses the generating function of Bernoulli numbers and the
  Fourier expansion $\sum_{n \neq 0} e^{2\pi i n x}/n^{2k}$.
\end{proof}

\subsection{Table of even zeta values}

\begin{table}[h]
  \centering
  \caption{Even zeta values $\zeta(2n)$ via Euler's formula.}
  \vspace{4pt}
  \begin{tabular}{@{}cclr@{}}
    \toprule
    $n$ & $\zeta(2n)$ (closed form) & Decimal approximation & $B_{2n}$ \\
    \midrule
    1 & $\dfrac{\pi^2}{6}$           & $1.6449340668\ldots$ & $\tfrac{1}{6}$     \\[6pt]
    2 & $\dfrac{\pi^4}{90}$          & $1.0823232337\ldots$ & $-\tfrac{1}{30}$   \\[6pt]
    3 & $\dfrac{\pi^6}{945}$         & $1.0173430620\ldots$ & $\tfrac{1}{42}$    \\[6pt]
    4 & $\dfrac{\pi^8}{9450}$        & $1.0040773562\ldots$ & $-\tfrac{1}{30}$   \\[6pt]
    5 & $\dfrac{\pi^{10}}{93555}$    & $1.0009945751\ldots$ & $\tfrac{5}{66}$    \\[6pt]
    6 & $\dfrac{691\pi^{12}}{638512875}$ & $1.0002460866\ldots$ & $-\tfrac{691}{2730}$ \\
    \bottomrule
  \end{tabular}
  \label{tab:even-zeta}
\end{table}

\subsection{Consequences}

Since every $B_{2n} \in \mathbb{Q}$, Euler's theorem immediately implies that
$\zeta(2n) \in \mathbb{Q} \cdot \pi^{2n}$ for all $n \geq 1$. Since $\pi$ is
transcendental (Lindemann 1882), each $\zeta(2n)$ is transcendental. Moreover,
the irrationality (and transcendence) of $\pi^{2n}$ was established well before
any result on odd zeta values.

The \emph{odd} case reveals the full depth of the problem: there is no known
analogue of the sine product expansion that would link $\zeta(2n+1)$ to $\pi$.

% ============================================================
\section{The Irrationality of $\zeta(3)$: Apéry's Theorem}
% ============================================================

\subsection{Statement}

\begin{theorem}[Apéry, 1979]
  The number $\zeta(3) = \sum_{n=1}^{\infty} n^{-3}$ is irrational.
\end{theorem}

This theorem was announced at the Journées Arithmétiques in Luminy in June~1978
and published in~1979. The result was initially met with scepticism---several
experts doubted the proof---but was verified within weeks by Henri Cohen and
others. Apéry received standing ovations at the announcement.

\subsection{The Apéry sequences}

The central idea of Apéry's proof is to construct two sequences of positive
integers $\{a_n\}$ and $\{b_n\}$ satisfying the same three-term linear recurrence,
converging to a rational approximation of $\zeta(3)$ so well that $\zeta(3)$
cannot be rational.

Define the Apéry numbers by
\[
  b_n = \sum_{k=0}^{n} \binom{n}{k}^2 \binom{n+k}{k}^2, \quad n \geq 0.
\]
The first few values are $b_0 = 1$, $b_1 = 5$, $b_2 = 73$, $b_3 = 1445$.

The numerator sequence $a_n$ is defined via
\[
  a_n = \sum_{k=0}^{n} \binom{n}{k}^2 \binom{n+k}{k}^2
        \sum_{j=1}^{k} \frac{(-1)^{j-1}}{j^3 \binom{j+n}{j}\binom{n}{j}},
\]
adjusted so that $a_n \in \mathbb{Z}$ after multiplication by a suitable
least common multiple.

\subsection{The key recurrence}

Both sequences satisfy the \emph{Apéry recurrence}:
\begin{equation}
  (n+1)^3 u_{n+1} - (34n^3 + 51n^2 + 27n + 5)\, u_n + n^3 u_{n-1} = 0.
  \label{eq:apery-recurrence}
\end{equation}
With initial conditions $(b_0, b_1) = (1, 5)$ this produces the $b_n$; with
$(a_0, a_1) = (0, 6)$ it produces the numerator sequence.

\subsection{Proof of irrationality}

The key estimates are:
\begin{enumerate}[label=(\roman*)]
  \item $b_n \zeta(3) - a_n \to 0$ as $n \to \infty$ (the ratio converges to $\zeta(3)$).
  \item The convergence rate is
        \[
          |b_n \zeta(3) - a_n| \sim C \cdot (\sqrt{2} - 1)^{4n}, \quad C > 0,
        \]
        so the convergence is super-exponentially fast.
  \item The denominator $b_n$ grows only polynomially (in fact $b_n \sim A \cdot
        (1 + \sqrt{2})^{4n} / n^3$ for some constant $A$).
\end{enumerate}

To prove irrationality: suppose $\zeta(3) = p/q$ for integers $p, q > 0$.
Then $|b_n \zeta(3) - a_n| = |b_n p/q - a_n| \geq 1/q$ for large $n$ (as a
non-zero integer divided by $q$), but the right-hand side tends to zero faster
than any polynomial, a contradiction.

\subsection{Beukers' integral formulation}

Beukers (1979) gave an elegant reformulation: define
\begin{align*}
  I_n &= \int_0^1 \!\int_0^1 \frac{-\log(xy)}{1 - xy}\, P_n(x)\, P_n(y)\,
         \mathrm{d}x\, \mathrm{d}y, \\
  J_n &= \int_0^1 \!\int_0^1 \!\int_0^1
         \frac{x^n(1-x)^n\, y^n(1-y)^n\, z^n(1-z)^n}
              {(1 - (1-z)xy)^{n+1}}\,
         \mathrm{d}x\, \mathrm{d}y\, \mathrm{d}z,
\end{align*}
where $P_n$ is the $n$-th Legendre polynomial. These integrals yield
$J_n = b_n \zeta(3) - a_n$ and provide the positivity and rate estimates needed.

\subsection{Irrationality measure}

The best known upper bound on the irrationality measure of $\zeta(3)$ is due to
Rhin and Viola (2001):
\[
  \mu(\zeta(3)) \leq 5.5138905...,
\]
meaning $|\zeta(3) - p/q| < q^{-5.5138905-\varepsilon}$ has only finitely many
solutions for any $\varepsilon > 0$.

% ============================================================
\section{Apéry's Constant $\zeta(3)$}
% ============================================================

\subsection{Numerical value and continued fraction}

The constant $\zeta(3)$ is now known as \emph{Apéry's constant}:
\[
  \zeta(3) = 1.2020569031595942854523691932819321312725\ldots
\]
Its continued fraction expansion begins
\[
  \zeta(3) = [1; 4, 1, 18, 1, 1, 1, 4, 1, 9, 9, 2, 1, 1, 1, 2, \ldots],
\]
exhibiting no obvious pattern. This is consistent with its conjectured (but
unproved) transcendence.

\subsection{Alternative series representations}

Several faster-converging series are known:

\begin{itemize}[leftmargin=2em]
  \item \textbf{Euler's formula:}
    \[
      \zeta(3) = \frac{5}{2} \sum_{n=1}^{\infty} \frac{(-1)^{n+1}}{n^3 \binom{2n}{n}}.
    \]
  \item \textbf{Alternating Dirichlet series:}
    \[
      \zeta(3) = \frac{4}{3} \sum_{n=1}^{\infty} \frac{(-1)^{n+1}}{n^3}.
    \]
  \item \textbf{Borwein--Crandall--Fee formula:}
    \[
      \zeta(3) = \frac{\pi^2}{7} \left( 1 - 4 \sum_{k=1}^{\infty}
      \frac{\zeta(2k)}{(2k+1)(2k+2)\, 4^k} \right).
    \]
\end{itemize}

\subsection{"Apéry's miracle"}

The term \emph{Apéry's miracle} refers to the mysterious fact that the
coefficients $b_n$ in Apéry's proof are integers---a priori by no means obvious
from the combinatorial sums. The integrality follows from the divisibility
properties of binomial coefficients, but there is no simple conceptual explanation.
Alfred van der Poorten's account ``A proof that Euler missed'' (1979) made Apéry's
argument accessible to a wide audience.

\subsection{Open status}

\begin{conjecture}
  $\zeta(3)$ is transcendental (i.e., not algebraic over $\mathbb{Q}$).
\end{conjecture}

This remains completely open (as of~2026). Proving irrationality was difficult
enough; transcendence requires entirely new methods.

% ============================================================
\section{Irrationality Among $\zeta(5), \zeta(7), \ldots$}
% ============================================================

\subsection{The Ball--Rivoal theorem}

\begin{theorem}[Ball--Rivoal, 2001]
  The $\mathbb{Q}$-vector space dimension of
  \[
    \operatorname{span}_{\mathbb{Q}}\{1,\, \zeta(3),\, \zeta(5),\, \ldots,\, \zeta(2n+1)\}
  \]
  is at least
  \[
    \frac{(1 + o(1))\,\ln n}{1 + \ln 2}
    \quad \text{as } n \to \infty.
  \]
  In particular, infinitely many of the values $\zeta(2k+1)$ are irrational.
\end{theorem}

\begin{proof}[Proof strategy]
  The proof constructs, for each large $r$, a linear form
  \[
    L_r = c_0^{(r)} + c_3^{(r)}\,\zeta(3) + c_5^{(r)}\,\zeta(5) + \cdots
          + c_{2r-1}^{(r)}\,\zeta(2r-1)
  \]
  with integer coefficients $c_j^{(r)}$, where $L_r \neq 0$ but $|L_r| \to 0$.
  The coefficients are obtained from a hypergeometric integral, and their
  smallness relative to the denominators forces the conclusion that the
  corresponding set $\{1, \zeta(3), \ldots, \zeta(2r-1)\}$ must contain at least
  $\sim \ln(r)/(1 + \ln 2)$ linearly independent elements over $\mathbb{Q}$.
\end{proof}

\subsection{The growth rate}

More precisely, the dimension satisfies
\[
  \dim_{\mathbb{Q}} \operatorname{span}_{\mathbb{Q}}\{1, \zeta(3), \zeta(5),
  \ldots, \zeta(2n+1)\} \geq \frac{\ln n}{1 + \ln 2} \approx \frac{\ln n}{1.6931}
\]
for large $n$. The lower bound grows logarithmically, confirming infinitely many
irrationals but leaving open the possibility (unlikely but not disproved) that
only a sparse subset is irrational.

\subsection{Zudilin's theorem}

\begin{theorem}[Zudilin, 2004]
  At least one of the four values
  \[
    \zeta(5),\quad \zeta(7),\quad \zeta(9),\quad \zeta(11)
  \]
  is irrational.
\end{theorem}

\begin{proof}[Proof strategy]
  Zudilin constructed a linear form
  \[
    L = a_0 + a_5 \zeta(5) + a_7 \zeta(7) + a_9 \zeta(9) + a_{11} \zeta(11),
    \quad a_i \in \mathbb{Z},
  \]
  with $L \neq 0$ but $|L| < 1$. If all four values were rational with a common
  denominator $D$, then $D \cdot L$ would be a non-zero integer strictly between
  $-1$ and $1$---a contradiction.

  The construction uses hypergeometric functions ${}_{7}F_6$ and the Siegel lemma
  to find suitable integer coefficients with controlled size.
\end{proof}

\begin{remark}
  Zudilin's theorem does \emph{not} identify which of the four values is irrational.
  The irrationality of $\zeta(5)$ individually remains an open conjecture.
\end{remark}

\subsection{Table of odd zeta values}

\begin{table}[h]
  \centering
  \caption{Odd zeta values and their arithmetic status.}
  \vspace{4pt}
  \begin{tabular}{@{}clll@{}}
    \toprule
    $n$ & $\zeta(n)$ (approximation) & Irrationality & Reference \\
    \midrule
    3  & $1.20205690315959\ldots$ & \textbf{Proved irrational} & Apéry (1979) \\
    5  & $1.03692775514337\ldots$ & Open conjecture           & --- \\
    7  & $1.00834927738192\ldots$ & Open conjecture           & --- \\
    9  & $1.00200839292608\ldots$ & Open conjecture           & --- \\
    11 & $1.00049418860409\ldots$ & Open conjecture           & --- \\
    13 & $1.00012271334757\ldots$ & Open conjecture           & --- \\
    15 & $1.00003058823630\ldots$ & Open conjecture           & --- \\
    \midrule
    \multicolumn{4}{l}{At least one of $\zeta(5), \zeta(7), \zeta(9), \zeta(11)$
    is irrational: \textbf{Proved} (Zudilin 2004)} \\
    \multicolumn{4}{l}{Infinitely many $\zeta(2k+1)$ are irrational: \textbf{Proved}
    (Ball--Rivoal 2001)} \\
    \bottomrule
  \end{tabular}
  \label{tab:odd-zeta}
\end{table}

% ============================================================
\section{Multiple Zeta Values}
% ============================================================

\subsection{Definition}

Multiple zeta values (MZVs) generalise the Riemann zeta function to multi-index
arguments. For integers $s_1 \geq 2$ and $s_2, \ldots, s_k \geq 1$, define
\[
  \zeta(s_1, s_2, \ldots, s_k) = \sum_{n_1 > n_2 > \cdots > n_k \geq 1}
  \frac{1}{n_1^{s_1} n_2^{s_2} \cdots n_k^{s_k}}.
\]
The integer $k$ is the \emph{depth} and $|s| = s_1 + \cdots + s_k$ is the
\emph{weight}.

\subsection{Examples and first values}

\begin{align*}
  \zeta(2) &= \frac{\pi^2}{6}, \\
  \zeta(3) &= 1.20205690\ldots \quad \text{(Apéry's constant)}, \\
  \zeta(2, 1) &= \zeta(3), \quad \text{(Euler's result)}, \\
  \zeta(1, 2) &= \zeta(3) \quad \text{(duality)}, \\
  \zeta(2, 3) &= 9\,\zeta(5)/2 - 2\,\zeta(2)\zeta(3).
\end{align*}

\subsection{The algebra of MZVs}

MZVs satisfy two families of algebraic relations:

\textbf{Shuffle product.} Integration-by-parts yields
\[
  \zeta(\mathbf{s}) \cdot \zeta(\mathbf{t}) = \sum_{\sigma \in \mathrm{Sh}(\mathbf{s}, \mathbf{t})}
  \zeta(\sigma),
\]
where the sum is over all shuffle permutations.

\textbf{Stuffle product (quasi-shuffle).} The Dirichlet series multiplication gives
a different identity:
\[
  \zeta(s_1) \cdot \zeta(s_2) = \zeta(s_1, s_2) + \zeta(s_2, s_1) + \zeta(s_1 + s_2).
\]

These two product structures impose strong linear relations among MZVs. The
\emph{double shuffle relations} (the two together) generate all known algebraic
relations, and it is conjectured (but not proved) that they generate \emph{all}
relations.

\subsection{Weight and period}

MZVs of weight $w$ conjecturally span a $\mathbb{Q}$-vector space of dimension
$d_w$ satisfying $d_w = d_{w-2} + d_{w-3}$ (Zagier's conjecture). The first
few values are:
\[
  d_1 = 0,\quad d_2 = 1,\quad d_3 = 1,\quad d_4 = 1,\quad d_5 = 2,
  \quad d_6 = 2, \quad d_7 = 3.
\]

% ============================================================
\section{Bernoulli Numbers and Euler--Maclaurin}
% ============================================================

\subsection{Definition and recurrence}

The Bernoulli numbers $B_n$ are defined by the generating function
\[
  \frac{t}{e^t - 1} = \sum_{n=0}^{\infty} B_n \frac{t^n}{n!}, \quad |t| < 2\pi.
\]
They satisfy the recurrence
\[
  \sum_{k=0}^{n} \binom{n+1}{k} B_k = 0, \quad n \geq 1,
\]
with $B_0 = 1$. One finds $B_1 = -1/2$, $B_{2k+1} = 0$ for $k \geq 1$, and

\begin{table}[h]
  \centering
  \caption{Bernoulli numbers $B_{2k}$.}
  \vspace{4pt}
  \begin{tabular}{@{}crr@{}}
    \toprule
    $k$ & $B_{2k}$ & Decimal \\
    \midrule
    1 & $1/6$               & $0.16\overline{6}$ \\
    2 & $-1/30$             & $-0.0\overline{3}$ \\
    3 & $1/42$              & $0.0238\ldots$ \\
    4 & $-1/30$             & $-0.0\overline{3}$ \\
    5 & $5/66$              & $0.07575\ldots$ \\
    6 & $-691/2730$         & $-0.25311\ldots$ \\
    7 & $7/6$               & $1.1\overline{6}$ \\
    \bottomrule
  \end{tabular}
\end{table}

\subsection{Euler--Maclaurin summation formula}

\begin{theorem}[Euler--Maclaurin]
  For a smooth function $f$ on $[1, N]$,
  \[
    \sum_{n=1}^{N} f(n) = \int_1^N f(x)\,\mathrm{d}x
    + \frac{f(1) + f(N)}{2}
    + \sum_{k=1}^{p} \frac{B_{2k}}{(2k)!}\bigl(f^{(2k-1)}(N) - f^{(2k-1)}(1)\bigr)
    + R_p,
  \]
  where $R_p$ is a remainder involving $f^{(2p+1)}$.
\end{theorem}

Applied to $f(x) = x^{-s}$, the Euler--Maclaurin formula yields the analytic
continuation of $\zeta(s)$ and, for negative integers, the special values
$\zeta(-n) = -B_{n+1}/(n+1)$.

\subsection{Kummer's congruences}

For a prime $p$ and $p \nmid k$, Kummer's congruences state
\[
  \frac{B_k}{k} \equiv \frac{B_{k'}}{k'} \pmod{p}
\]
whenever $k \equiv k' \pmod{(p-1)}$ and $p \nmid k$.
These are the cornerstone of $p$-adic interpolation of zeta values.

% ============================================================
\section{$p$-adic Zeta Functions}
% ============================================================

\subsection{Motivation}

The rational values $\zeta(1-n) = -B_n/n$ for positive integers $n$ (odd $n$,
with appropriate sign conventions) satisfy Kummer's congruences, which suggests
that $n \mapsto \zeta(1-n)$ can be continuously extended to $p$-adic integers.

\subsection{Kubota--Leopoldt $p$-adic $L$-functions}

\begin{theorem}[Kubota--Leopoldt, 1964]
  For every prime $p$ and Dirichlet character $\chi$ with $\chi(-1) = (-1)^k$,
  there exists a unique continuous $p$-adic function
  \[
    L_p(s, \chi): \mathbb{Z}_p \to \mathbb{C}_p
  \]
  such that, for all positive integers $k$ with $k \equiv 0 \pmod{p-1}$ (or
  $k$ even in the case $\chi = \mathbf{1}$),
  \[
    L_p(1-k, \chi) = -\frac{B_{k,\chi}}{k} \cdot \left(1 - \chi(p) p^{k-1}\right),
  \]
  where $B_{k,\chi}$ is the generalised Bernoulli number.
\end{theorem}

\subsection{$p$-adic interpolation and Iwasawa theory}

The Kubota--Leopoldt function is the central object of Iwasawa theory. In the
cyclotomic $\mathbb{Z}_p$-extension $\mathbb{Q}_\infty$ of $\mathbb{Q}$, the
projective limit of class groups yields a finitely generated $\Lambda$-module
(where $\Lambda = \mathbb{Z}_p[[T]]$ is the Iwasawa algebra), whose characteristic
polynomial is closely related to $L_p(s, \chi)$.

The \emph{Main Conjecture of Iwasawa Theory} (proved by Mazur--Wiles in~1984 for
$\mathbb{Q}$, using modular curves) asserts that the characteristic ideal of this
module equals the ideal generated by the $p$-adic $L$-function.

\subsection{Connection to odd zeta values}

The $p$-adic zeta function interpolates the values $\zeta(1-k)$ for even positive
$k$, not the odd values directly. The lack of a $p$-adic interpolation for odd
positive integers is itself a manifestation of the arithmetic mystery surrounding
$\zeta(2k+1)$.

% ============================================================
\section{Connections to Modular Forms}
% ============================================================

\subsection{Eisenstein series and $\zeta(2k)$}

The Eisenstein series of weight $2k$ for $\mathrm{SL}_2(\mathbb{Z})$ is
\[
  G_{2k}(\tau) = \sum_{\substack{(m,n) \in \mathbb{Z}^2 \\ (m,n) \neq (0,0)}}
  \frac{1}{(m\tau + n)^{2k}}, \quad k \geq 2, \quad \tau \in \mathbb{H}.
\]
Its Fourier expansion is
\[
  G_{2k}(\tau) = 2\zeta(2k) + \frac{2(2\pi i)^{2k}}{(2k-1)!}
  \sum_{n=1}^{\infty} \sigma_{2k-1}(n)\, q^n, \quad q = e^{2\pi i \tau},
\]
where $\sigma_r(n) = \sum_{d|n} d^r$. This directly expresses $\zeta(2k)$ as the
constant term of an Eisenstein series, explaining the appearance of $\pi^{2k}$.

\subsection{Why odd values are different}

There is no Eisenstein series of odd weight for $\mathrm{SL}_2(\mathbb{Z})$
(since $\mathrm{SL}_2(\mathbb{Z})$ contains $-I$, which forces modular forms
of odd weight to vanish). This is a deep structural reason why $\zeta(2k+1)$
does not have a simple expression in terms of $\pi$.

\subsection{The Ramanujan--Berndt formula for $\zeta(2k+1)$}

Ramanujan (and later Berndt) proved the formula
\[
  \zeta(2k+1) = \frac{2^{2k} \pi^{2k+1}}{(2k+1)!}
  \left[
    - B_{2k+1} + 2 \sum_{n=1}^{\infty} \frac{\sigma_{2k+1}^*(n)}{e^{2\pi n} - 1}
  \right]
\]
where $\sigma_r^*$ involves divisor sums and $B_{2k+1}$ are generalised Bernoulli
numbers, but this formula requires the ``correction'' from Eisenstein-like series.
It does \emph{not} reduce $\zeta(2k+1)$ to a rational multiple of $\pi^{2k+1}$,
confirming the arithmetic difference.

\subsection{Period integrals}

In the language of motives, even zeta values $\zeta(2n)$ are \emph{periods of
mixed Tate motives} that split over $\mathbb{Q}$, whereas odd zeta values are
related to the \emph{motivic cohomology} of $\mathrm{Spec}(\mathbb{Z})$. This
motivic perspective, developed by Beilinson, Deligne, and Goncharov, provides
the deepest framework for understanding the arithmetic of zeta values.

% ============================================================
\section{Open Problems}
% ============================================================

\begin{conjecture}[Transcendence of $\zeta(3)$]
  The constant $\zeta(3) = 1.2020569\ldots$ is transcendental.
\end{conjecture}

\begin{conjecture}[Irrationality of $\zeta(5)$]
  The value $\zeta(5) = 1.0369277\ldots$ is irrational.
\end{conjecture}

\begin{conjecture}[Irrationality of all odd zeta values]
  For every $k \geq 1$, the value $\zeta(2k+1)$ is irrational.
\end{conjecture}

\begin{conjecture}[$\zeta(2n+1)/\pi^{2n+1}$ irrational]
  For every $k \geq 1$, the ratio $\zeta(2k+1)/\pi^{2k+1}$ is irrational (and
  in fact not algebraic).
\end{conjecture}

\begin{conjecture}[Algebraic independence]
  The numbers $\pi, \zeta(3), \zeta(5), \zeta(7), \ldots$ are algebraically
  independent over $\mathbb{Q}$.
\end{conjecture}

\subsection*{Status summary}

\begin{table}[h]
  \centering
  \caption{Open problems concerning odd zeta values (as of 2026).}
  \vspace{4pt}
  \begin{tabular}{@{}lll@{}}
    \toprule
    Claim & Status & Reference \\
    \midrule
    $\zeta(3)$ irrational & \textbf{PROVED} & Apéry (1979) \\
    $\zeta(2k+1)$ irrational ($\infty$ many) & \textbf{PROVED} & Ball--Rivoal (2001) \\
    At least one of $\zeta(5),\ldots,\zeta(11)$ irrational & \textbf{PROVED} & Zudilin (2004) \\
    $\zeta(3)$ transcendental & OPEN conjecture & --- \\
    $\zeta(5)$ irrational & OPEN conjecture & --- \\
    All $\zeta(2k+1)$ irrational & OPEN conjecture & --- \\
    $\zeta(2k+1)/\pi^{2k+1}$ irrational & OPEN conjecture & --- \\
    Algebraic independence of $\pi, \zeta(3), \zeta(5), \ldots$ & OPEN conjecture & --- \\
    \bottomrule
  \end{tabular}
\end{table}

\subsection*{Approaches and outlook}

Substantial progress would likely require one or more of the following:
\begin{enumerate}[leftmargin=2em]
  \item A better understanding of the \emph{motivic Galois group} and its action
        on periods.
  \item New hypergeometric identities yielding smaller linear forms in $\zeta(5)$.
  \item A $p$-adic interpolation of odd zeta values (which would require new
        structures beyond Kubota--Leopoldt).
  \item Effective versions of the Lindemann--Weierstrass theorem for polylogarithms.
\end{enumerate}

% ============================================================
\section{References}
% ============================================================

\begin{thebibliography}{99}

\bibitem{Apery1979}
R.~Apéry,
\textit{Irrationalité de $\zeta(2)$ et $\zeta(3)$},
Astérisque \textbf{61} (1979), 11--13.

\bibitem{BallRivoal2001}
K.~Ball and T.~Rivoal,
\textit{Irrationalité d'une infinité de valeurs de la fonction zeta aux entiers impairs},
Inventiones Mathematicae \textbf{146} (2001), 193--207.

\bibitem{Beukers1979}
F.~Beukers,
\textit{A note on the irrationality of $\zeta(2)$ and $\zeta(3)$},
Bulletin of the London Mathematical Society \textbf{11} (1979), 268--272.

\bibitem{Euler1740}
L.~Euler,
\textit{De summis serierum reciprocarum},
Commentarii Academiae Scientiarum Petropolitanae \textbf{7} (1740), 123--134.

\bibitem{RhinViola2001}
G.~Rhin and C.~Viola,
\textit{The group structure for $\zeta(3)$},
Acta Arithmetica \textbf{97} (2001), 269--293.

\bibitem{Rivoal2000}
T.~Rivoal,
\textit{La fonction zêta de Riemann prend une infinité de valeurs irrationnelles
aux entiers impairs},
Comptes Rendus de l'Académie des Sciences (Paris), Série I \textbf{331} (2000),
267--270.

\bibitem{Zudilin2004}
W.~Zudilin,
\textit{Arithmetic of linear forms involving odd zeta values},
Journal de Théorie des Nombres de Bordeaux \textbf{16} (2004), 251--291.

\bibitem{vanderPoorten1979}
A.~van der Poorten,
\textit{A proof that Euler missed \ldots\ Apéry's proof of the irrationality of
$\zeta(3)$},
The Mathematical Intelligencer \textbf{1} (1979), 195--203.

\bibitem{Zagier1994}
D.~Zagier,
\textit{Values of zeta functions and their applications},
in: \textit{First European Congress of Mathematics}, Vol.~II,
Birkhäuser, Basel, 1994, pp.~497--512.

\bibitem{KubotaLeopoldt1964}
T.~Kubota and H.~Leopoldt,
\textit{Eine $p$-adische Theorie der Zetawerte I},
Journal für die reine und angewandte Mathematik \textbf{214/215} (1964), 328--339.

\end{thebibliography}

\end{document}
